\section{The precise continuation}
This tranche completes a fixed-target divisor problem in the original Type-II
cofactor box. It does not assert occupancy at every prime. Starting from an
actual exterior state with a negative-four-integer shape, it gives the entire
same-grade middle return, including every noncanonical divisor and its inverse.
The noncanonical repair has a sharp finite bound depending only on the shape
coefficient. That bound is realized by original exterior states at infinitely
many hard primes. Conversely, for every fixed shape coefficient $t\ge4$, an
explicit reduced prime progression gives original exterior states whose entire
same-grade middle box is empty. Their known exterior solutions remain present.

The Type-I/II coordinate identities and the square-divisor cofactor criterion
are antecedents \cite[\S2, Propositions 2.2 and 2.6]{ET}. The preceding source
\cite[negative-four-square transfer]{Parent} supplied $t=c^2$ when
$c\mid(h+1)/4$. Here the arithmetic restriction is neither omitted nor made into
an unspecified admissibility hypothesis: it is solved exactly as a divisor
problem. No historical-priority determination is claimed.

\section{Original coordinates and a general negative shape}
Let $p\equiv1\pmod4$ be prime. An original exterior state consists of
\[
 p/4<a<p/2,\quad R=4a-p,\quad u\mid a^2,\quad R\mid4u+1.
\]
Retain $a=\prod_l l^{e_l}$, the full box $-e_l\le\beta_l\le e_l$, and
$u=\prod_l l^{e_l+\beta_l}$. Set
\[
 g=\gcd(a,u),\quad (h,r,s)=(g^2/u,u/g,a/g),\quad
 \kappa=(pr+s)/R.
\]
Then $a=hrs$, $u=hr^2$, $\gcd(r,s)=1$, and the ordered denominators are
\[
 (x,y,z)=(a,hs\kappa,phr\kappa),\qquad
 \frac4p=\frac1x+\frac1y+\frac1z.                         \tag{1}
\]
Multiplication proves the identity, and its ordered inverse is
$u=a^2/(Ry-pa)$. All normalization factors are positive integers. Let
\[
 v=hs^2=a^2/u,\quad D=(4u+1)/R=(r+\kappa)/s,
 \quad \alpha=v-R,\quad\beta=v-D.
\]
The symbol $\beta$ here is not an exponent $\beta_l$. The original identities are
\[
 RD=4hr^2+1,\quad pD+1=4hr\kappa,\quad r+\kappa=sD,
\]
\[
 h\mid\alpha\beta-1,\quad p\equiv\alpha\pmod h,
 \quad p\beta\equiv1\pmod h.                             \tag{2}
\]
For every original divisor $d\mid a^2$, reciprocity gives $(d/R)=(d/p)$.
For an odd prime $l\mid a$, use $R\equiv-p\pmod l$ and $R\equiv3\pmod4$;
for 2 use the supplementary law. The exterior target $u\equiv-1/4\pmod R$
therefore gives $(h/p)=(u/p)=-1$. These are the classical laws, used in their
original factor coordinates \cite[\S6.4]{Hatcher}.

Suppose one ordered shape coordinate is $-4t$, $t\ge1$. Since $R,D\equiv3\pmod4$,
\[
 h\equiv3\pmod4,\qquad s\ \hbox{odd},\qquad
 \gcd(h,t)=1,\qquad (t/h)=1.                             \tag{3}
\]
Indeed $hs^2$ is 3 modulo four. Equation (2) makes $h$ coprime to $4t$.
For $\alpha$, reciprocity gives $(h/p)=(p/h)=(-4t/h)=-(t/h)$; the inverse
congruence gives the same result for $\beta$. Also $p>2h$ because $a\ge h$.

Put $j=(h+1)/4$ and define the full original target set
\[
 \boxed{\mathcal A_t(j)=\{W>0:W\mid j^2,\ W\equiv t\pmod{4j-1}\}.}       \tag{4}
\]
There is no primality condition on $h=4j-1$ in (4).

\begin{theorem}[Complete same-grade channel return]\label{thm:return}
For every original exterior state as above and every $W\in\mathcal A_t(j)$, set
\[
 U=\begin{cases}W,&\alpha=-4t,\\j^2/W,&\beta=-4t,\end{cases}
 \quad Q=h,\quad R_M=(p+4U)/h,\quad a_M=(p+R_M)/4.          \tag{5}
\]
These are exactly the original cofactor-$h$ middle choices, with the indicated
ordered shape branch retained. In particular $U\mid j^2$, $0<R_M<p$, and
\[
 g_M=\gcd(j,U),\quad
 h_M=g_M^2/U,\quad r_M=U/g_M,\quad\lambda_M=j/g_M,
 \quad s_M=R_M\lambda_M-r_M                              \tag{6}
\]
returns the original middle state and its ordered denominators
\[
 (a_M,ph_Ms_M\lambda_M,ph_Mr_M\lambda_M).
\]
\end{theorem}
\begin{proof}
In the $\alpha$ case $p\equiv-4t\pmod h$, so the cofactor congruence is exactly
$W\equiv t$. In the $\beta$ case $p\equiv-(4t)^{-1}\pmod h$ and
$j^2\equiv16^{-1}\pmod h$, so $U=j^2/W$ satisfies $p+4U\equiv0$ precisely
when $W\equiv t$. Every divisor involved is a unit modulo $h$.
Since $U\le j^2$, $h\ge3$ and $p>2h$,
\[
 4U\le(h+1)^2/4<(h-1)p.
\]
Thus $R_M<p$, and its positive integer value is 3 modulo four. Prime valuations
in $U\mid j^2$ give $h_Mr_M\lambda_M=j$, $h_Mr_M^2=U$, and
$\gcd(r_M,\lambda_M)=1$. Directly
\[
 a_M=R_Mj-U=h_Mr_Ms_M>0.
\]
A common prime of $r_M$ and $R_M$ would divide $p$ by $hR_M=p+4U$, whereas
$r_M\le j<p$. Hence $\gcd(r_M,s_M)=1$. These identities prove the middle
normalization, the first-half range, and the reciprocal identity. Conversely a
cofactor-$h$ middle word has $U\mid j^2$ and the same gate; taking $W=U$ or
$j^2/U$ recovers (4). Neither $U<a_M$ nor $r_M<s_M$ is imposed: the ordered
orientation is kept, or complemented explicitly afterward.
\end{proof}

\section{An exact disjoint classification of every available divisor}
The next theorem is pure integer arithmetic. In particular, $t$ need not be a
square and $h$ need not be prime.

\begin{theorem}[Fixed-target capacity completion]\label{thm:capacity}
Let $t,j\ge1$ and $h=4j-1>t$. The set $\mathcal A_t(j)$ is the disjoint union of:
\begin{enumerate}
\item the canonical word $W=t$, present exactly when $t\mid j^2$;
\item the companion word $W=4tj$, present exactly when $4t\mid j$;
\item the finite words described by
\[
 1\le\ell\le\left\lfloor\frac{t-2}{4}\right\rfloor,
 \quad w\mid\ell^2,\quad
 n=\frac{t-w}{4\ell+1}\in\mathbb Z_{>0},                  \tag{7}
\]
\[
 V=\ell^2/w,\quad j=\ell+4nV,\quad h=4j-1>t,
 \quad W=j^2/V=t+nh.                                    \tag{8}
\]
\end{enumerate}
In the last item the displayed $j$ must equal the input $j$. Each word occurs
once. The inverse, for $W\ne t$, is
\[
 n=(W-t)/h,\quad V=j^2/W,\quad\ell=j-4nV.                \tag{9}
\]
For $\ell=0$ it is the companion. Otherwise $w=\ell^2/V$ is the unique label
in (7).
\end{theorem}
\begin{proof}
If $W\ne t$, then $W=t+nh$ for some $n\ge1$, because $0<t<h$. Let $V=j^2/W$.
Reduction of $16WV=16j^2=(h+1)^2$ modulo $h$ gives
\[
 16tV=1+kh
\]
for an integer $k\ge1$. Modulo four, $k\equiv1$, so $k=4\ell+1$ with
$\ell\ge0$. Comparing the two exact products gives
\[
 h+2=k+16nV,\qquad j=\ell+4nV,
 \qquad t=(4\ell+1)n+\frac{\ell^2}{V}.                  \tag{10}
\]
If $\ell=0$, then $n=t$, $j=4tV$, and $W=4tj$. This is equivalent to
$4t\mid j$ and coexists with the canonical divisor. If $\ell>0$, (10) proves
$V\mid\ell^2$, and $w=\ell^2/V\ge1$. Therefore $t\ge4\ell+2$, giving (7).
All labels are recovered by (9). Conversely (7)--(8) give
\[
 j^2/V=w+8\ell n+16n^2V=t+n(4j-1),
\]
an integer divisor of $j^2$ in the required class. The three cases are disjoint
by $W=t$, $\ell=0$, or $\ell>0$.
\end{proof}

The small-grade domain $h\le t$ is separately finite:
$j\le\lfloor(t+1)/4\rfloor$. Literal enumeration there is retained; the theorem
is not asserted on that domain. A reduced residue $t\bmod h$ may be used to test
a word, but it must not replace the original shape coefficient in a bound.

\begin{theorem}[Sharp finite-repair bound]\label{thm:bound}
Every finite word in Theorem~\ref{thm:capacity} satisfies
\[
 \boxed{h\le t^2-3t+1.}                                \tag{11}
\]
Equality holds exactly when $w=n=1$ and $t=4\ell+2$. In particular, if the
canonical divisor is absent but $\mathcal A_t(j)$ is nonempty, (11) holds.
For $h>\max(t,t^2-3t+1)$ the exact count is
\[
 \boxed{|\mathcal A_t(j)|=
  \mathbf1_{t\mid j^2}+\mathbf1_{4t\mid j}.}             \tag{12}
\]
\end{theorem}
\begin{proof}
Use $n=(t-w)/(4\ell+1)$ to write
\[
 h+1=4\ell+\frac{16\ell^2(t-w)}{(4\ell+1)w}
 \le4\ell+\frac{16\ell^2(t-1)}{4\ell+1}.
\]
The last expression is strictly increasing in real $\ell\ge0$. Since
$\ell\le(t-2)/4$, it is at most
$(t-2)+(t-2)^2=t^2-3t+2$. Equality requires both $w=1$ and
$\ell=(t-2)/4$, and then $n=1$. The companion cannot repair a missing canonical
word because $4t\mid j$ implies $t\mid j^2$. The classification proves (12).
\end{proof}

For every $\ell\ge1$ the arithmetic equality family is
\[
 t=4\ell+2,\quad j=\ell(4\ell+1),\quad h=4j-1,
 \quad W=(4\ell+1)^2,\quad V=\ell^2.                    \tag{13}
\]
Here $t\nmid j^2$: the odd number $2\ell+1>1$ is coprime to both factors of
$j$. Thus (13) is a genuine repair, not a second name for the canonical word.
Not every finite word is a repair: at
\[
 t=289,\quad j=1224,\quad h=4895,
 \qquad\mathcal A_t(j)=\{289,5184\},
\]
the canonical word and a finite word coexist. The latter has
$(n,V,\ell,w)=(1,289,68,16)$. Its coefficient is two, even though its occupied
residue is a single class.

The complete finite search uses at most
\[
 \sum_{\ell\le(t-2)/4}\tau_{\rm div}(\ell^2)
 \le L(1+\log L)^2,\quad L=\lfloor(t-2)/4\rfloor\ge1.   \tag{14}
\]
Indeed $\tau_{\rm div}(\ell^2)\le d_3(\ell)$ primewise, and
$\sum_{abc\le L}1\le L(\sum_{a\le L}1/a)^2$.  This counts candidate labels
under unit-cost integer arithmetic; it is not a bit-complexity estimate that
includes integer factorization.  It is a finite classification in the original
integers, not an assertion that a successful word always exists.

There is also a strict smaller-parameter return on each finite word. Put
$g_0=\gcd(\ell,V)$ and
\[
 H_0=g_0^2/V,\quad r_0=V/g_0,\quad \lambda_0=\ell/g_0.
\]
Then $\ell=H_0r_0\lambda_0$, $V=H_0r_0^2$, and $\gcd(r_0,\lambda_0)=1$.
The middle normalization in (6) is
\[
 \begin{array}{c|ccc}
 &h_M&r_M&\lambda_M\\\hline
 \alpha&H_0&\lambda_0+4nr_0&r_0\\
 \beta&H_0&r_0&\lambda_0+4nr_0.
 \end{array}                                           \tag{15}
\]
In particular $h_M\le\ell\le(t-2)/4$. This decreases a parameter of the
finite capacity reduction; it is not a well-founded proof of initial ES occupancy.

\section{Entire shape lines with forced original availability}
\begin{theorem}[Automatic negative-shape returns]\label{thm:auto}
Let $p\equiv1\pmod{24}$ be prime and start from an original exterior state.
The canonical word is always available on the three lines
\[
 \alpha=-4t,\qquad t\in\{1,2,3\},
\]
and on the nine lines
\[
 \boxed{\beta=-4t,\qquad t\mid36.}                      \tag{16}
\]
Every such line has the original middle return (5)--(6), with no bound on the
other shape coordinate. The coefficient $t$ need not be a square.
\end{theorem}
\begin{proof}
For $\alpha$, (3) gives $(2/h)=1$ at $t=2$; together with $h\equiv3\pmod4$,
this forces $h\equiv7\pmod8$, hence $2\mid j$. At $t=3$, the supplementary
quadratic law gives $h\equiv11\pmod{12}$ and $3\mid j$. The case $t=1$ is literal.

For $\beta$, use the stronger actual equations, not just their Jacobi projection.
Here $D=hs^2+4t$ and $r+\kappa=sD$ is odd. Thus $r\kappa$ is even. Since
$pD+1=4hr\kappa$ and $p\equiv1\pmod8$, necessarily $D\equiv7\pmod8$, whence
\[
 h\equiv7-4t\pmod8.
\]
If $2\mid t$, this gives $2\mid j$.
If $3\mid t$, all equalities in the rest of this paragraph are modulo three.
The case $h=0$ makes $D=0$ and contradicts
$pD+1=4hr\kappa$. If $h=1$ and $s=0$, the same equation would give
$1=-r^2$, impossible. If $h=1$ and $s=\pm1$, then $D=1$, $\kappa=s-r$, and
$2=r(s-r)$; the right side has only values zero and one. Thus $h=2\pmod3$,
so $3\mid j$. Every prime exponent of $t\mid36$ is at most two. Therefore
these two forced divisibilities imply $t\mid j^2$, and Theorem~\ref{thm:return}
applies.
\end{proof}

For the canonical $W=t$, the squarefree part of $h_M$ is exactly that of $t$.
This follows from $h_M=g_M^2/U$ and, in the inverse branch, $U=j^2/t$.
The preceding square-target theorem is the special case in which this part is one.
We do not claim that (16) lists every possible automatic $\beta$ coefficient.
Actual complete failures at $t=5,8,27$ are given below.

\section{The quadratic bound is sharp on actual hard-prime sources}
\begin{theorem}[Prime realization of the finite-repair boundary]\label{thm:sharp}
For every $t=6+420z$, $z\ge0$, infinitely many primes $p\equiv1\pmod{840}$ have
an original exterior state with $\alpha=-4t$,
\[
 h=t^2-3t+1,\quad R=t^2+t+1,\quad s=1,
\]
whose canonical cofactor word is unavailable but whose finite word in (13)
returns an original middle state. Thus the bound (11) is attained within the
original prime arithmetic, for unbounded $t$.
\end{theorem}
\begin{proof}
Put $\ell=(t-2)/4$, $j=\ell(4\ell+1)$, so $h=4j-1$ and the word is
$W=(t-1)^2$. The two integers $h,R$ are coprime, since their difference is $4t$
and $h\equiv R\equiv1\pmod t$. For the displayed progression of $t$,
$\gcd(hR,210)=1$. Also $t^3\equiv1\pmod R$. Consequently
\[
 r_0=(t+2R)/4\in\mathbb Z,
 \qquad16tr_0^2\equiv1\pmod R.
\]
Since $h\equiv-4t\pmod R$, the latter is exactly
$R\mid4hr_0^2+1$.
Choose $r$ by the explicit compatible CRT conditions
\[
 r\equiv r_0\pmod R,\qquad
 r\equiv\frac{R+1}{4}h^{-1}\pmod{210}.
\]
Then $p=4hr-R$ varies in a progression of modulus $840hR$, is 1 modulo 840,
and is coprime to the entire modulus: modulo $h$ it is $-R$ and modulo $R$
it is $4hr_0$. Dirichlet's theorem supplies infinitely many prime members
\cite[Theorem 18.1]{Dirichlet}. For all sufficiently large positive members,
$a=hr$, $u=hr^2$ is an original first-half exterior state; $R<2hr$ verifies the
 only upper range issue. Its shape is $h-R=-4t$. Formula (13) proves canonical
 unavailability and supplies the sharp finite repair. The target is reconstructed by
Theorem~\ref{thm:return}. Prime values are supplied by a reduced \emph{linear}
progression; no assertion about primes represented by a higher-degree polynomial
has been inserted.
\end{proof}

A small explicit original member, in another hard class, is
\[
 \boxed{p=825241,\quad
 (a,R,u,h,r,s,\kappa)=(206321,43,2240439739,19,10859,1,208402140).}    \tag{17}
\]
Its coefficient is $t=6$, $j=5$, and the entire capacity box is $\{1,5,25\}$.
Only $25\equiv6\pmod{19}$ succeeds. The original middle return is
\[
 (a_M,R_M,U,h_M,r_M,s_M,\lambda_M)
 =(217170,43439,25,1,5,43434,1),                           \tag{18}
\]
with ordered denominators $(217170,825241\cdot43434,825241\cdot5)$.
This is a nonsquare coefficient and a genuinely noncanonical return.

\section{Every other alpha coefficient has infinite original same-grade obstructions}
\begin{lemma}[Exact limitation of the character-only budget]
For $t\ge1$, let the product run over every prime, including \(2\):
\[
 K(t)=\prod_{l^e\parallel t}l^{\lceil e/2\rceil},\qquad M=4K(t).
\]
For odd \(A\) coprime to \(t\), write
\(\chi_t(A)=(t/A)\) for the quadratic Kronecker--Jacobi character determined
by the squarefree part of \(t\); its conductor divides \(M\).  Among units
\(A\pmod M\) with \(A\equiv3\pmod4\) and \(\chi_t(A)=1\), the implication
$A\equiv-1\pmod M$ holds for all such $A$ exactly when $t=1,2,3$.
\end{lemma}
\begin{proof}
The Jacobi character is well-defined modulo $M$ by reciprocity, including the
factor at two. It takes value one at $-1$. If $t$ is a square, it is trivial,
and the indicated set has $\varphi(M)/2$ elements. If $t$ is nonsquare, its
restriction to the units 1 modulo four is nontrivial: at an odd prime appearing
to odd exponent choose a nonsquare by CRT, with the other coordinates one; if
only the factor two has odd exponent, choose the coordinate 5 modulo eight.
Thus it is independent of the mod-four sign character, and the indicated set
has $\varphi(M)/4$ elements. A multiple of four with totient two is $M=4$;
those with totient four are $8,12$. The square/non-square alternatives give
exactly $t=1,2,3$. In each singleton case the member is $-1$.
\end{proof}

\begin{theorem}[An actual obstruction for every fixed $t\ge4$]\label{thm:obstruct}
For every integer $t\ge4$, there are infinitely many hard primes possessing an
original exterior state with $\alpha=-4t$ but no middle state at cofactor equal
to that state's original grade $h$, even after the \emph{entire} divisor box
$\operatorname{Div}(j^2)$ is tested. These primes already have the constructed
exterior solution. This is not an Erd\H{o}s--Straus counterexample.
\end{theorem}
\begin{proof}
Choose a unit class $A\pmod{4K(t)}$ from the preceding lemma other than $-1$.
Then $A\equiv3\pmod4$ and $(t/A)=1$. Lift it by CRT to a reduced class modulo
$\operatorname{lcm}(4K(t),840)$ such that, for $l=3,5,7$, both $R$ and
$R-4t$ are units modulo $l$. If $l\mid K(t)$ this is already true, because
$l\mid t$ and $A$ is a unit. Otherwise one need only avoid $0,4t\pmod l$,
which leaves a choice even for $l=3$. Choose an odd lift modulo eight as well.
Dirichlet's theorem gives a prime $R$ in this class as large as required.
Set $h=R-4t$ and choose it so that
\[
 h>\max(4t,t^2-3t+1).
\]
Then $h\equiv3\pmod4$, $\gcd(h,210)=\gcd(h,R)=1$, and
$h\not\equiv-1\pmod{4K(t)}$. Thus $t\nmid j^2$ for $j=(h+1)/4$.
   Since \(\chi_t(R)=(t/R)=1\), choose \(q\) with \(q^2\equiv t\pmod R\)
   and set \(r_0\equiv(4q)^{-1}\pmod R\).  Then
\[
 16tr_0^2\equiv1\pmod R.
 \]
 Solve
\[
 r\equiv r_0\pmod R,\qquad
 r\equiv\frac{R+1}{4}h^{-1}\pmod{210}.
\]
Again $p=4hr-R$ is a reduced progression modulo $840hR$, with $p\equiv1\pmod{840}$.
Dirichlet gives infinitely many prime members. Take $a=hr$, $u=hr^2$, $s=1$.
Then $R\mid4u+1$, $\gcd(p,a)=1$, and $p>2a$ follows already from $h>4t$ and
$r\ge1$. The original shape is $\alpha=h-R=-4t$.
Theorem~\ref{thm:bound} now gives $\mathcal A_t(j)=\varnothing$: neither the
canonical word nor its companion is available and the finite-exception range
has been passed. Theorem~\ref{thm:return} identifies this with failure of the
complete cofactor-$h$ middle source.
\end{proof}

For $t=4$ one explicit construction is
\[
 R=59,\quad h=43,\quad j=11,\quad
 r=8415+12390n,\quad
 \boxed{p=1447321+2131080n.}                              \tag{19}
\]
The progression is reduced and 1 modulo 840. At $n=3$, $p=7840561$ is prime,
with
\[
 (a,R,u,h,r,s,\kappa)
 =(1960155,59,89353665675,43,45585,1,6057830054).
\]
The three capacity divisors are $1,11,121$, of residues $1,11,35$ modulo 43,
not four. Every prime in (19) nevertheless has its displayed exterior solution.

The beta automatic theorem has genuine limits too. The following are original
hard-prime sources, and complete enumeration of their small $j^2$ boxes proves
failure at their old grade:
\[
\begin{array}{c|r|r|r|r|r|r|r}
 p&a&u&h&r&s&\kappa&t\ (\beta=-4t)\\\hline
346201&158815&43864703&23&1381&5&1654&8\\
593401&188100&39710000&11&1900&9&7091&27\\
2727841&929071&464321099&11&6497&13&17930&5.
\end{array}                                             \tag{20}
\]
The case $h\le t$ in the middle row is checked literally, not attributed to
Theorem~\ref{thm:capacity}. No infinite beta obstruction theorem is asserted.

\section{Original inverse fibres and a necessary prose correction}
For fixed $p,h,t$ on an $\alpha$ branch, put $N=(p+4t)/h$. Its complete original
fibre is
\[
 s\mid N,\quad s+N/s\equiv0\pmod4,\quad r=(s+N/s)/4,
 \quad a=hrs,\quad u=hr^2,
\]
\[
 \boxed{\gcd(r,s)=1,\quad p/4<a<p/2,\quad
        R=4a-p\mid4hr^2+1.}                             \tag{21}
\]
The shape equation is $4rs-s^2=N$; the extra gate in (21) is necessary and not
implied by this factorization. For a beta branch, put $F=(4tp+1)/h$ and enumerate
\[
 1\le s\le\lfloor(p-1)/(2h)\rfloor,\quad
 D=hs^2+4t,\quad\Delta_s=s^2(D^2-p)-F.
\]
Keep the positive squares $k^2=\Delta_s$ with $sD-k$ positive even, and set
$r=(sD-k)/2$, $\kappa=sD-r$, retaining the original coprimality, range and gate.
These are equivalences from $4hr(sD-r)=pD+1$ and $k=\kappa-r>0$.

The predecessor's prose alpha-fibre formula omitted the final gate in (21).
Its executable inverse-fibre enumeration did test that gate. At
$p=5209,h=95,t=4$, the displayed divisor equation has $N=55$. The choice
$s=5,r=4$ satisfies its factorization, coprimality, shape and first-half conditions,
but gives
\[
 a=1900,\quad R=2391,\quad u=1520,\qquad
 (4u+1)\bmod R=1299.
\]
It is not an exterior state. This correction changes the textual claimed inverse
domain, not the predecessor's forward theorem or its executable fibre filter.
The old source is preserved, and \texttt{CORRECTION.md} identifies the affected
passage and file hash.

For the new maps, retain $(p,\text{original E state},\text{side},t,j,W)$ and its
canonical/companion/finite label. The M output recovers $Q=h$ and $j$; (21) and
its beta counterpart recover precisely the possible E sources once $t$ and the
side remain. The finite labels are inverted by (9). Across two ordered branches,
alpha word $W_1$ and beta word $W_2$ give the same M word exactly when
$W_1W_2=j^2$. Multiple original sources above that word are not identified.

On the free abelian group of these labelled sources the coefficient pushforward
sums each actual target fibre. Its kernel is generated by same-target differences.
No division by a fibre size is used to prove positivity. Reducing a positive even
coefficient modulo two can yield a supported zero; it does not remove its
original words. The two-word example following (13) is an explicit control.

\section{Verification scope and the remaining milestone}
The verifier exhausts (4) for $t\le120$, $j\le5000$, $4j-1>t$, and compares every
integer word with Theorem~\ref{thm:capacity}. It separately verifies all finite
profiles and one hundred members of the sharp arithmetic family, the original
prime-realizable progressions, and the character-only classification through 120.
It also enumerates every original first-half E/M divisor at primes
$p\equiv1\pmod4$ through 3000, retaining actual factorization and channel records,
and checks the returned original states. The fixed larger examples use complete
trial-division primality checks. Their empty cofactor boxes are exhausted.

The independent program imports neither the main verifier nor predecessor code.
It uses literal divisors of $j^2$, coprime primitive $r,s$ sources, primewise
Euler-symbol evaluation, and a complementary-divisor parameter for the finite
profiles. Normal and optimized command outcomes, hashes and scopes belong to the
execution receipt. Multiple runs are not independent mathematical review.

The completed unbounded statements here concern the \emph{fixed-target capacity
problem} and the original shape lines in Theorem~\ref{thm:auto}. The obstruction
in Theorem~\ref{thm:obstruct} proves why increasing the old-grade divisor search
cannot repair every negative-square alpha source. It does not exclude a different
original middle cofactor or an exterior solution, which the construction already
supplies. No initial state has been forced at an arbitrary hard prime. Turn~7's
universal existence milestone remains outstanding; no Turn~8 witness guarantee,
Lean build, new global verification range or priority claim follows.
